% ─────────────────────────────────────────────────────────────────────────────
% CAPÍTULO 3 — CI/CD, Frontend y Aplicación Android
% ─────────────────────────────────────────────────────────────────────────────
\chapter{CI/CD, Despliegue y App Android}

\section{Pipeline CI/CD — GitHub Actions}

El repositorio incluye un workflow completo en
\code{.github/workflows/deploy.yml} que cubre dos etapas: validación de
calidad y despliegue zero-downtime al VPS.

\begin{figure}[h]
\centering
\begin{tikzpicture}[
  node distance=0.5cm and 1.0cm,
  box/.style={draw=hermesdark,fill=hermeslight,rounded corners=5pt,
              minimum width=3.2cm,minimum height=1.0cm,
              font=\sffamily\small,align=center},
  good/.style={draw=greenok,fill=greenok!10,rounded corners=5pt,
               minimum width=3.2cm,minimum height=1.0cm,
               font=\sffamily\small,align=center},
  bad/.style={draw=redalert,fill=redalert!10,rounded corners=5pt,
              minimum width=3.2cm,minimum height=1.0cm,
              font=\sffamily\small,align=center},
  arr/.style={->,>=Stealth,thick,hermesblue},
  rarr/.style={->,>=Stealth,thick,redalert,dashed}
]
\node[box]  (push)   {Git push\\a \texttt{main}};
\node[box,right=of push]   (lint)   {Etapa 1\\ruff · black · mypy};
\node[good,right=of lint]  (deploy) {Etapa 2\\SSH → VPS\\docker compose up};
\node[good,right=of deploy](health) {Health check\\LiteLLM + Hermes};
\node[good,below=1cm of health](ok) {\textbf{ÉXITO}\\Producción actualizada};
\node[bad, below=1cm of deploy](rb) {\textbf{ROLLBACK}\\git checkout HEAD\textasciitilde1\\docker compose up};

\draw[arr]  (push)   -- (lint);
\draw[arr]  (lint)   -- node[above,font=\tiny\sffamily]{pass} (deploy);
\draw[arr]  (deploy) -- (health);
\draw[arr]  (health) -- node[right,font=\tiny\sffamily]{OK} (ok);
\draw[rarr] (health) -- node[left,font=\tiny\sffamily\color{redalert}]{FAIL} (rb);
\draw[rarr] (lint)   |- node[below,font=\tiny\sffamily\color{redalert}]{fail} ++(0,-1.8) -- ($(push)+(0,-1.8)$);
\end{tikzpicture}
\caption{Pipeline CI/CD del Hermes Stack}
\end{figure}

\subsection{Etapa 1 — Lint \& Validation}

Se ejecuta en todo PR y push. Nunca accede al servidor de producción.

\begin{lstlisting}[language=yaml,caption=Jobs de validación en deploy.yml]
lint-and-test:
  runs-on: ubuntu-24.04
  steps:
    - uses: actions/checkout@v4
    - uses: actions/setup-python@v5
      with:
        python-version: "3.12"
        cache: pip
    - run: pip install ruff black mypy aiohttp websockets tenacity prometheus_client
    - run: ruff check ./hermes
    - run: black --check ./hermes
    - run: mypy ./hermes --ignore-missing-imports
\end{lstlisting}

\subsection{Etapa 2 — Despliegue al VPS}

Solo se ejecuta en push a \code{main} y requiere que Etapa 1 pase.

\begin{lstlisting}[language=yaml,caption=Despliegue SSH con rollback automático]
deploy:
  needs: [lint-and-test]
  if: github.ref == 'refs/heads/main' && github.event_name == 'push'
  steps:
    - name: Deploy via SSH
      run: |
        ssh $VPS_USER@$VPS_HOST "
          cd /root
          git pull origin main
          docker compose up -d --build
          sleep 15

          # Health check autenticado LiteLLM
          (source .env && curl -f \
            -H 'Authorization: Bearer $LITELLM_MASTER_KEY' \
            http://127.0.0.1:4000/health) \
          || (git checkout HEAD~1 && docker compose up -d --build && exit 1)

          # Health check Hermes Agent
          curl -f http://127.0.0.1:8080/health \
          || (git checkout HEAD~1 && docker compose up -d --build && exit 1)
        "
\end{lstlisting}

\subsection{Secretos requeridos en GitHub}

\begin{table}[h]
\centering
\caption{GitHub Actions Secrets (Settings → Secrets and variables → Actions)}
\begin{tabularx}{\textwidth}{lX}
\toprule
\textbf{Secret}    & \textbf{Valor} \\
\midrule
\envvar{VPS\_HOST}    & IP pública del servidor (\code{172.236.x.x}) \\
\envvar{VPS\_USER}    & Usuario SSH (\code{root}) \\
\envvar{VPS\_SSH\_KEY}& Clave privada SSH (su pública en \code{/root/.ssh/authorized\_keys}) \\
\envvar{VPS\_PORT}    & Puerto SSH (default \code{22}) \\
\bottomrule
\end{tabularx}
\end{table}

\begin{table}[h]
\centering
\caption{Variables en \texttt{/root/.env} del VPS (nunca salen del servidor)}
\begin{tabularx}{\textwidth}{lX}
\toprule
\textbf{Variable}               & \textbf{Descripción} \\
\midrule
\envvar{LITELLM\_MASTER\_KEY}   & Clave maestra del proxy LiteLLM \\
\envvar{OPENROUTER\_API\_KEY}   & API key de OpenRouter \\
\envvar{GEMINI\_API\_KEY}       & API key de Google Gemini \\
\envvar{ELEVENLABS\_API\_KEY}   & API key de ElevenLabs \\
\envvar{ELEVENLABS\_VOICE\_ID}  & ID de voz seleccionada en ElevenLabs \\
\bottomrule
\end{tabularx}
\end{table}

\subsection{Watchdog del Docker Daemon}

\code{docker-watchdog.sh} es un servicio systemd que supervisa el daemon
de Docker y los contenedores críticos cada 60 segundos, independientemente
de Autoheal:

\begin{lstlisting}[language=bash,caption=Lógica central del watchdog]
for item in "litellm-router:litellm" "hermes-agent:hermes" "whisper-stt:whisper-stt"; do
  container="${item%%:*}"
  service="${item#*:}"
  status=$(docker inspect --format='{{.State.Status}}' "$container" 2>/dev/null \
           || echo "missing")
  if [ "$status" != "running" ]; then
    log "Contenedor $container caido. Levantando servicio $service..."
    cd "$STACK_DIR" && docker compose up -d --no-deps "$service"
  fi
done
\end{lstlisting}

\section{Configuración Caddy y TLS Automático}

\service{Caddy} gestiona los certificados TLS mediante el protocolo ACME
(Let's Encrypt) automáticamente. Solo requiere que los puertos 80 y 443
estén abiertos en el firewall.

\begin{lstlisting}[caption=/etc/caddy/Caddyfile completo]
{
    email erikjosehernandez@gmail.com
}

hermes.el80.space    { reverse_proxy 127.0.0.1:8080 }
litellm.el80.space   { reverse_proxy 127.0.0.1:4000 }
grafana.el80.space   { reverse_proxy 127.0.0.1:3000 }
docs.el80.space      { reverse_proxy 127.0.0.1:3001 }
syncthing.el80.space { reverse_proxy 127.0.0.1:8384 }
files.el80.space     { reverse_proxy 127.0.0.1:8095 }
\end{lstlisting}

\begin{lstlisting}[language=bash,caption=Instalación de Caddy (setup-caddy.sh)]
# Repositorio oficial
curl -1sLf 'https://dl.cloudsmith.io/public/caddy/stable/gpg.key' \
  | gpg --dearmor -o /usr/share/keyrings/caddy-stable-archive-keyring.gpg

# Firewall
ufw allow 80/tcp && ufw allow 443/tcp

# Habilitar y arrancar
systemctl enable caddy && systemctl restart caddy
\end{lstlisting}

\section{Blueprint de App Nativa Android}

La hoja de ruta incluye una app Android nativa que se conecta directamente
al Hermes Stack como cliente de voz y control.

\subsection{Stack tecnológico recomendado}

\begin{table}[h]
\centering
\caption{Stack tecnológico de la app Android}
\begin{tabularx}{\textwidth}{lX}
\toprule
\textbf{Capa} & \textbf{Tecnología} \\
\midrule
Lenguaje      & Kotlin + Coroutines + Flow \\
UI Framework  & Jetpack Compose (Material 3) \\
Audio         & MediaRecorder / AudioRecord → PCM 16 kHz mono \\
Red           & OkHttp WebSocket para streaming de voz \\
LLM directo   & Google AI Studio SDK (Gemini Flash en el device) \\
Persistencia  & Room DB para historial de conversaciones \\
DI            & Hilt \\
\bottomrule
\end{tabularx}
\end{table}

\subsection{Secuencia de voz en la app Android}

\begin{figure}[h]
\centering
\begin{tikzpicture}[
  node distance=0.4cm and 0.8cm,
  comp/.style={draw=hermesdark,fill=hermeslight,rounded corners=4pt,
               minimum width=2.8cm,minimum height=0.8cm,
               font=\sffamily\scriptsize,align=center},
  arr/.style={->,>=Stealth,thick,hermesblue}
]
\node[comp] (mic)  {Micrófono\\Android};
\node[comp,right=of mic]  (buf)  {Buffer PCM\\16 kHz};
\node[comp,right=of buf]  (ws)   {WebSocket\\hermes.el80.space};
\node[comp,right=of ws]   (hrm)  {Hermes\\Agent};
\node[comp,right=of hrm]  (ll)   {LiteLLM\\→ LLM};
\node[comp,below=1cm of ll](tts) {ElevenLabs\\TTS WS};
\node[comp,left=of tts]   (sp)   {Speaker\\Android};

\draw[arr] (mic) -- (buf);
\draw[arr] (buf) -- (ws);
\draw[arr] (ws)  -- (hrm);
\draw[arr] (hrm) -- (ll);
\draw[arr] (ll)  |- (tts);
\draw[arr] (tts) -- (sp);
\end{tikzpicture}
\caption{Flujo de voz en la app Android}
\end{figure}

\subsection{Roadmap mobile}

\begin{enumerate}[label=\textbf{Fase \arabic*:}]
  \item Autenticación con \envvar{LITELLM\_MASTER\_KEY} vía HTTPS
  \item WebSocket de voz bidireccional con Hermes Agent
  \item Dashboard de health (pull de \domain{docs.el80.space/api/health})
  \item Llamadas REST directas a LiteLLM para chat sin voz
  \item Integración de Gemini on-device vía Google AI Studio SDK
\end{enumerate}
